% Unnumbered post29 diagnostic. Not promoted or assembled.
\subsection*{Canonical stable overorders with unbounded cohomology}

This note tests what follows from the canonical-order conductor
comparison alone. It supplies an actual canonical full order, stable
anchors, and Azumaya overorders. The cohomology is unbounded even
when the divisorial defect is fixed and the relevant saturation
quotient is zero. The geometric Brauer period is one and the anchor
rank is two. Thus this is \emph{not} a counterexample to the sharper
rank-equals-period construction or to the existence of a good anchor.
It asserts no failure of a representation-count theorem.

\paragraph{One fixed curve, bundle and cone.}
Work over $k=\overline{\mathbf F}_3$. Choose a smooth projective
curve $C$ of genus $g\ge2$ and a stable rank-two bundle $G$ with
$\det G\simeq\mathcal O_C$ and degree zero. Such a choice can be
made concretely; an elementary construction is given below.
Stability gives
\[
 H^0(C,\mathcal End G)=k\,1.
\]
Since two is invertible, the trace decomposition is
\[
 \mathcal End G=\mathcal O_C\,1\oplus H,
 \qquad H=\ker(\operatorname{tr}),\quad \operatorname{rk}H=3,
 \quad H^0(C,H)=0.
\]
The trace pairing on $H$ is perfect, and $\det H\simeq\mathcal O_C$.
Choose a sufficiently positive normally generated ample line bundle
$L$ such that $H\otimes L$ is generated by global sections and
$H^0(C,H\otimes L^{-1})=0$. Serre generation and, by duality,
Serre vanishing provide such a fixed choice. Write $d=\deg L>0$.
Multiplication by a nonzero section of a positive power of $L$
then shows $H^0(C,H\otimes L^{-a})=0$ for all $a\ge1$.

For a fully explicit normal-cone setup, one may take
$C\subset\mathbf P^2_k$ to be the Fermat quartic
$X^4+Y^4+Z^4=0$ and $L=\mathcal O_C(e)$ for a sufficiently large
$e$ prime to three. This curve is smooth because its partial
derivatives have no simultaneous projective zero, and $g=3$.
Its homogeneous coordinate ring is a normal standard graded
two-dimensional domain: it is a hypersurface and is regular
off the vertex. The section ring for $L$ is its $e$-th Veronese
subring, which is normal and generated in degree one.

Let $R$ be the completed local ring of this fixed cone at its
vertex, and let $Y\to\operatorname{Spec}R$ be the base change of
the standard cone resolution
$Y^{\rm alg}=\operatorname{Tot}(L^{-1})$. Its exceptional curve
is $E\simeq C$, with
\[
 \mathcal O_E(-E)=L,\qquad E^2=-d.
\]
The excellent normal completion and regularity of the proper
base-changed resolution are the same already checked facts used
in the earlier cone packets. All subsequent choices vary only
the bundle, not $R,Y,E,G,$ or $L$.

\paragraph{The order is genuinely the canonical full order.}
Let $F$ be the pullback of $G$ to $Y$, and let $U=Y\setminus E$.
Put
\[
 P_0=j_*\mathcal End(F|_U),\qquad
 P=\operatorname{full}_Y(P_0),
 \quad j:U\hookrightarrow\operatorname{Spec}R.
\]
The analogous graded section module before completion is
\[
 \bigoplus_{a\in\mathbf Z}
       H^0(C,\mathcal End(G)\otimes L^a).
\]
The scalar part is the section ring. The trace-zero part has
no terms of degree at most zero, and its degree-one sections
generate $H\otimes L$. On the total-space resolution the
image of these degree-one sections is exactly
$(\operatorname{pr}^*H)(-E)$: the fibre-linear factor vanishes
once along the zero section. Every higher-degree section lies
in the same lattice. Thus the actual torsion-free image of
the pullback is already locally free and equals
\[
       P=\mathcal O_Y\,1\oplus(\operatorname{pr}^*H)(-E).
                                                        \qquad(1)
\]
Flat completion preserves this description, the finite section
module and its reflexive pullback. In particular (1) is not an
arbitrarily chosen smaller order. It is the canonical full
pullback of the reflexive punctured endomorphism algebra.
It is a unital subalgebra of $\mathcal End(F)$: products of
two trace-zero summands vanish at least twice along $E$.
Its action on $F|_E=G$ is scalar.

\paragraph{An unbounded family of stable anchors for this same order.}
For every sufficiently large integer $n$ choose a surjection
\[
                       G\twoheadrightarrow L^n.

\]
Such a surjection exists: $G^\vee\otimes L^n$ is globally
generated of rank two. For each point the sections vanishing
there form a codimension-two subspace, so the incidence of a
point and a vanishing section has dimension at most one less
than the section space. Properness of $C$ then supplies a
nowhere-zero section, equivalently the stated quotient.
Let $i:E\hookrightarrow Y$ and define
\[
              K_n=\ker(F\longrightarrow i_*L^n).
                                                        \qquad(2)
\]
Locally this is the elementary kernel of a free rank-two
module surjecting onto a line on a Cartier divisor, and is
therefore a vector bundle of rank two. It restricts to the
same $F|_U$ for every $n$, and
\[
                \det K_n=\det F\otimes\mathcal O_Y(-E)

\]
is independent of $n$.

The scalar part of (1) preserves (2). Its remaining part
sends $F$ into $F(-E)\subset K_n$, so it also preserves
$K_n$. Thus every $K_n$ is an actual $P$-stable anchor.
Set
\[
                         B_n=\mathcal End(K_n).

\]
This is a locally free Azumaya algebra of rank four, with
the same specified punctured algebra as $P$. Trace integrality
on $K_n$ gives the actual inclusions
\[
          P\subset B_n\subset P^\vee\subset P(2E).
                                                        \qquad(3)
\]
For the last inclusion one can compute directly from (1):
\[
 P^\vee=\mathcal O_Y\,1\oplus(\operatorname{pr}^*H)(E),
 \qquad
 P(2E)=\mathcal O_Y(2E)\,1\oplus(\operatorname{pr}^*H)(E).

\]
The pairing on the scalar summand is multiplication by two,
so there is no characteristic-three degeneracy. The fixed
Noether Cartier cycle may also be enlarged once to dominate
$2E$; its bound is independent of $n$.

All these orders have $\det B_n\simeq\mathcal O_Y$.
The determinant divisor of $P\hookrightarrow B_n$ is exactly
$3E$, independently of $n$: it is supported on $E$, whereas
$\det P=\mathcal O_Y(-3E)$; intersection with $E$, whose
self-intersection is nonzero, fixes its coefficient.
Moreover $\Gamma(Y,B_n)=P_0$, since $P\subset B_n$ and
the canonical $P$ already has the entire punctured section
module. These strengthened comparisons still do not yield
the cohomology bound below.

\paragraph{Cohomology grows at every fixed twist.}
Restricting (2) to $E$ retains its nonzero Tor term and gives
\[
       0\longrightarrow L^{n+1}\longrightarrow K_n|_E
              \longrightarrow L^{-n}\longrightarrow0.
                                                        \qquad(4)
\]
Here the last line is the kernel of $G\to L^n$, since
$\det G=\mathcal O_C$. Composing the quotient and inclusion
in (4) embeds a line into the endomorphism bundle:
\[
             L^{2n+1}\hookrightarrow\mathcal End(K_n|_E).

\]
For every fixed integer $c$ this gives
\[
 h^0\bigl(E,B_n(-cE)|_E\bigr)
     \ge (2n+1+c)d+1-g.

\]
The rank-four endomorphism bundle has degree zero, so
Riemann--Roch implies
\[
 h^1\bigl(E,B_n(-cE)|_E\bigr)
      \ge (2n+1-3c)d+3(g-1).
                                                        \qquad(5)
\]
The right side tends to infinity. For each fixed positive
integer $b$, restriction from $bE$ onto $E$ is surjective
on $H^1$, because the kernel is a coherent sheaf on a proper
one-dimensional scheme and has no $H^2$. Likewise restriction
from $Y$ onto $E$ is surjective on $H^1$, using the established
coherent $H^2$ vanishing for this proper map with fibres of
dimension at most one over an affine base. Hence the same
unboundedness holds on every fixed thick cycle and on $Y$.
In particular no fixed multiple of the scalar comparison
cycle gives the proposed order-valued vanishing for all
these $P$-stable choices.

\paragraph{The finite-support saturation defect can already be zero.}
Take the punctured target object to be the punctured anchor
itself. Then its canonical Hom pullback is $Q=P$, the induced
object is $K_n$, and the corresponding coarse Hom module is
$E_n=B_n$. Therefore
\[
                   Q\cdot B_n=B_n=E_n.

\]
There is no finite-support saturation quotient at all, and
$E_n$ is free of rank one over $B_n$. Its background Hom
cohomology still grows as in (5). Thus bounding only the
saturation quotient or its generator number would not,
by itself, prove the desired quantitative presentation.

\paragraph{Elementary existence of the initial stable bundle.}
Choose a point $P_C\in C$. The extension space
\[
 V=\operatorname{Ext}^1(\mathcal O_C(P_C),
                         \mathcal O_C(-P_C))
   =H^1(C,\mathcal O_C(-2P_C))

\]
has dimension $g+1$. For each $Q_C\in C$, the kernel of
\[
 H^1(\mathcal O_C(-2P_C))\longrightarrow
 H^1(\mathcal O_C(-2P_C+Q_C))

\]
is a line, by the point exact sequence and the vanishing
of both negative-degree $H^0$ groups. These lines vary
algebraically: dually they are the evaluations of the
globally generated line $\omega_C(2P_C)$. Their union,
including zero, has dimension at most two. Since $g+1\ge3$,
choose an extension class outside that union and let $G$
be the resulting extension of $\mathcal O_C(P_C)$ by
$\mathcal O_C(-P_C)$.

If a line subbundle of $G$ had nonnegative degree, its
map to $\mathcal O_C(P_C)$ would be nonzero, hence would
identify it with $\mathcal O_C(P_C-D)$ for an effective
divisor of degree at most one. For $D=0$ the extension
would split. For $D=Q_C$ its pullback would split, placing
the chosen class in one of the excluded lines. Both are
impossible. Thus $G$ is stable of degree zero. Stability
implies that every endomorphism is scalar: a nonzero proper
image and kernel would contradict the strict degree-zero
slope inequality, and subtracting an eigenvalue over $k$
then proves the assertion. This proves all the initial
bundle hypotheses without a new finite-cover theorem.

\paragraph{Exact scope.}
The example is over one fixed complete normal surface in
characteristic three, with one fixed reflexive punctured
algebra and its canonical full order. It imposes actual
order stability, trace compatibility, a fixed Cartier
conductor, fixed determinant and fixed divisorial defect.
The geometric Brauer class is nevertheless trivial, while
the anchor rank is two. The minimal period-one anchor has
rank one and is not this example. No conclusion is drawn
against the rank-$M$ optimized construction, an appropriately
normalized distinguished anchor, or the requested hereditary
count. No Frobenius periodicity of $G$ is asserted.

The calculations use only the already checked curve
Riemann--Roch/duality, Serre generation, normal-cone
completion, regular-surface elementary-kernel facts and
the canonical order comparison. This is a separate
unnumbered diagnostic, not a four-audited manuscript unit.
